% =================================================
% 4. AGENT TAXONOMY  (included from main.tex with % =================================================
% 4. AGENT TAXONOMY  (included from main.tex with % =================================================
% 4. AGENT TAXONOMY  (included from main.tex with % =================================================
% 4. AGENT TAXONOMY  (included from main.tex with \input{section4_taxonomy})
% Needs: tikz, graphicx, booktabs, tabularx, array (already in the preamble).
% =================================================
\DeclareRobustCommand{\pfull}{\tikz[baseline=-0.6ex]\fill (0,0) circle (0.8ex);}
\DeclareRobustCommand{\phalf}{\tikz[baseline=-0.6ex]{\draw (0,0) circle (0.8ex);
  \fill (0,-0.8ex) arc (-90:90:0.8ex) -- cycle;}}
\DeclareRobustCommand{\pnone}{\tikz[baseline=-0.6ex]\draw (0,0) circle (0.8ex);}
\providecommand{\rot}[1]{\rotatebox{75}{#1}}

\section{Agent Taxonomy}
\label{sec:taxonomy}

In this section we analyse the agents of Section~3 with the theory seen in the course:
agent properties, agent architectures and types of agents by purpose
\cite{lecture2, lecture3, wooldridge2009}. To keep the tables short we use these
abbreviations: Flight Data Assistant (FDA), Trajectory Planner (TP), Weather Data
Collector (WDC), Weather Hazard Analyst (WHA), Radar Controller (RC), Conflict Analyst
(CA), Executive Controller (EC), Verification Officer (VO), ATC Supervisor (SUP) and ATC
Report Officer (RO).

% -------------------------------------------------
\subsection{Agent Properties}
\label{sec:properties}

Table~\ref{tab:properties} shows which of the ten properties of Lecture~3 each agent has.
When we filled it in, we saw that three features of our design explain most of the
results. First, the system runs only once and then stops. Second, all the calculations
are done by deterministic tools. Third, the agents exchange information through
intermediate files with a fixed structure (Pydantic models).

\begin{table}[H]
\centering
\small
\caption{Agent properties. \pfull{} = yes, \phalf{} = partially, \pnone{} = no.}
\label{tab:properties}
\begin{tabular}{@{}l*{10}{c}@{}}
\toprule
\textbf{Agent} & \rot{Flexibility} & \rot{Reactivity} & \rot{Proactiveness} & \rot{Social ability}
  & \rot{Rationality} & \rot{Reasoning} & \rot{Learning} & \rot{Autonomy}
  & \rot{Temporal continuity} & \rot{Mobility} \\
\midrule
FDA & \phalf & \phalf & \pnone & \phalf & \pfull & \pnone & \pnone & \phalf & \pnone & \pnone \\
TP  & \pnone & \phalf & \pnone & \phalf & \pfull & \pnone & \pnone & \phalf & \pnone & \pnone \\
WDC & \pnone & \phalf & \phalf & \phalf & \pfull & \pnone & \pnone & \phalf & \pnone & \pnone \\
WHA & \phalf & \phalf & \phalf & \phalf & \pfull & \phalf & \pnone & \phalf & \pnone & \pnone \\
RC  & \pnone & \phalf & \pnone & \phalf & \pfull & \pnone & \pnone & \phalf & \pnone & \pnone \\
CA  & \phalf & \phalf & \pfull & \phalf & \pfull & \pfull & \pnone & \phalf & \pnone & \pnone \\
EC  & \pfull & \pfull & \pfull & \pfull & \pfull & \pfull & \phalf & \pfull & \pnone & \pnone \\
VO  & \pnone & \pfull & \pnone & \phalf & \pfull & \pnone & \pnone & \phalf & \pnone & \pnone \\
SUP & \pfull & \pfull & \pfull & \pfull & \pfull & \pfull & \pnone & \pfull & \pnone & \pnone \\
RO  & \phalf & \phalf & \phalf & \phalf & \pfull & \phalf & \pnone & \phalf & \pnone & \pnone \\
\bottomrule
\end{tabular}
\end{table}

The data and detection agents (FDA, TP, WDC, RC, VO) have few properties, because they only
apply a fixed procedure. On the other hand, EC and SUP have almost all of them, since they
deal with the open part of the problem: deciding how to solve each conflict.

% -------------------------------------------------
\subsection{Justification of Non-obvious Property Classifications}
\label{sec:properties-justification}

\textbf{Reactivity.} Our environment does not change while the system is running
(Section~2), so in the usual sense \cite{lecture3} our agents would not be very reactive.
We decided to look instead at the changes in the information that each agent receives
from the others. Most agents get their input only once, so we marked them as partially
reactive. EC, VO and SUP are different: in each iteration of the verification loop they
receive new flight plans, a rejection or new feedback, and they have to adapt to it.

\textbf{Proactiveness.} CA does more than pass the list of conflicts on. It decides which
ones are more urgent and gives the Tactical crew an order to work on them. RO is
partially proactive, because it points out the critical cases to the human controller
without being asked.

\textbf{Social ability.} Our agents do not use an agent communication language like
FIPA-ACL. They write and read intermediate files, which works in a similar way to a
blackboard system \cite{lecture4}, and the Pydantic models act as a shared ontology. The
only real two-way dialogue happens between EC and SUP: EC proposes a manoeuvre and SUP
accepts or rejects it, which is close to the propose, accept-proposal and reject-proposal
acts of FIPA \cite{fipa_acl}.

\textbf{Rationality.} A rational agent acts towards its goals as far as its beliefs allow
\cite{lecture3}. The problem is that a language model can have wrong beliefs, for example a
wrong distance between two aircraft. That is why all the numbers in our system come from
deterministic tools, and every manoeuvre is checked by the Safety Validation crew before
it reaches the report.

\textbf{Learning.} No agent learns over time: each execution is independent and nothing is
stored between runs. We only marked EC as partial, because inside one run it uses the
feedback of SUP to improve its next proposal. Saving past conflicts and the solutions that
were accepted (case-based reasoning) could be a future improvement.

\textbf{Autonomy.} In the course we saw that the user can limit the autonomy of an agent,
and that some decisions should not be fully delegated to a machine \cite{lecture3}. In ATC
this is very clear. EC is free to choose how to solve a conflict, but it cannot apply its
decision directly: SUP has to approve it, and conflicts that are still unsolved after two
attempts go to a human controller. In the forum we also agreed that the iteration counter
is kept in the flow state and not by the language model, so the agent cannot decide to
skip the limit.

\textbf{Temporal continuity and mobility.} Agents with temporal continuity are running all
the time, and they are not ``once-only computations'' that read an input, produce an
output and stop \cite{lecture3}. Our system works in this second way, so none of our
agents has this property, although a real ATC system would need it. None of them is mobile
either, since all of them run on the same computer.

% -------------------------------------------------
\subsection{Agent Architectures}
\label{sec:architectures}

We classified the agents as reactive, deliberative or hybrid \cite{lecture2}. Since all
CrewAI agents use a language model internally, the question for us was how each agent
makes its decisions: by following a fixed rule, or by reasoning about its goals.
Table~\ref{tab:classification} shows the result.

\begin{table}[H]
\centering
\small
\caption{Architecture and type of each agent.}
\label{tab:classification}
\begin{tabularx}{\textwidth}{@{}l l l X@{}}
\toprule
\textbf{Agent} & \textbf{Architecture} & \textbf{Type} & \textbf{Comment} \\
\midrule
Flight Data Assistant  & Reactive     & Information   & Combines GraphML and JSON sources \\
Trajectory Planner     & Reactive     & Information   & Builds the trajectories for the other crews \\
Weather Data Collector & Reactive     & Information   & Loads the weather snapshot \\
Weather Hazard Analyst & Hybrid       & Information   & Turns weather data into hazard zones \\
Radar Controller       & Reactive     & Information   & Runs the detection tool \\
Conflict Analyst       & Hybrid       & Collaborative & Prioritises conflicts for Tactical \\
Executive Controller   & Deliberative & Collaborative & BDI-style planning of manoeuvres \\
Verification Officer   & Reactive     & Collaborative & Same detection rule as the RC \\
ATC Supervisor         & Hybrid       & Collaborative & Coordinates with the EC \\
ATC Report Officer     & Deliberative & Interface     & Assists the human controller \\
\bottomrule
\end{tabularx}
\end{table}

\textbf{Reactive agents.} FDA, TP, WDC, RC and VO do not make plans. They just map a
situation to an action ($S \rightarrow A$) \cite{lecture2}. RC and VO are a good example,
because their main rule is ``if two aircraft are closer than 5~NM horizontally and
1000~ft vertically, report a conflict'' \cite{icao_doc4444}. It reminded us of the
obstacle-avoidance rule in the reactive layer of TouringMachines \cite{ferguson1992}. We
want this kind of agent here because the check runs in every iteration, so it has to be
fast and reliable.

\textbf{Deliberative agents.} EC can be described with the BDI model \cite{bratman1987}.
Its \emph{beliefs} are the conflicts, the hazard zones, the aircraft performance and the
feedback of SUP. Its \emph{desire} is to solve all the conflicts with as little delay or
extra distance as possible. The manoeuvres it chooses are its \emph{intentions}, and the
concrete actions (climb, reroute or delay) are its \emph{plans}. When SUP rejects a
manoeuvre, EC drops that intention and looks for another one, in line with the idea that
an agent keeps its intentions only while it believes they can be achieved
\cite{cohen1990}. We also see RO as deliberative, because it has to plan the report from
the results of all the crews.

\textbf{Hybrid agents.} WHA, CA and SUP mix a reactive part (thresholds and tools) with a
deliberative part (interpreting, prioritising or deciding). For SUP we found that the
three layers of InteRRaP \cite{muller1996} describe it quite well. Its behaviour layer
rejects any manoeuvre that creates a new conflict. Its planning layer looks for the
reason and suggests something else. Its cooperative layer sends structured feedback to EC.

If we look at the whole system, it also has a hybrid structure. The tools work as a
reactive layer, the language-model agents as a deliberative layer, and the CrewAI Flow
with its routers decides which one acts, a bit like the control framework of
TouringMachines. The routers are deterministic, so the most critical decisions never
depend on the language model.

% -------------------------------------------------
\subsection{Agent Types by Purpose}
\label{sec:types}

Agents can also be classified by purpose as collaborative, interface, information or
facilitator agents \cite{lecture3, nwana1996}. Our system has several of these types
working together.

\textbf{Collaborative agents.} Air traffic control is one of the examples of collaborative
agents given in the course \cite{lecture3}, and in our design CA, EC, VO and SUP solve the
conflicts together through proposals, checks and feedback. Since all the agents share the
same goal, they are benevolent. Using the cooperation hierarchy of Lecture~5
\cite{doran1997}, we would place our system among the cooperative MAS with explicit
communication between deliberative agents. It also follows the steps of distributed
problem solving \cite{lecture4}. The problem is split into crews at design time, each crew
solves one part, the results are shared through the intermediate files, and the Reporting
crew puts them together in the final report.

\textbf{Information agents.} These agents collect, filter and combine data from different
sources \cite{lecture3}. FDA merges three input files that come in different formats, WDC
loads the weather data, and WHA turns it into hazard zones. TP and RC then give the
decision agents a clear view of the trajectories and the conflicts. In the forum we chose
to load a stored weather file instead of downloading live data, so WDC always works with
the same snapshot and the results can be repeated.

\textbf{Interface agent.} RO is the only agent that works for a person, the controller who
reads the report. We call it an interface agent only in a limited sense, because interface
agents usually learn from their user \cite{lecture3} and ours does not.

\textbf{Facilitators.} There is no facilitator agent in our design. The CrewAI Flow takes
this role: it decides which crew runs next and passes the data between them. We think a
facilitator agent would only make sense if new agents could join the system while it is
running.

% -------------------------------------------------
\subsection{Deterministic Conflict Tool and Agentification}
\label{sec:agentification}

The conflict detection code from the course notebook is not an agent, so we had to decide
how to connect it to the MAS. The course presents three options: a translator, a wrapper,
or rewriting the code as an agent \cite{lecture3, genesereth1994}.

A translator was not a good option, because the notebook calculates the validation errors
but does not return them, and its output has no structure. Rewriting it as an agent would
mean letting the language model calculate the distances, and we did not want that in a
safety-critical system. In the end we chose a \textbf{wrapper}. We moved the code into a
Python module with small changes (Pydantic models and returning the validation errors)
and added a CrewAI tool on top of it \cite{crewai_docs}. The algorithm is the same, and our
tests check that it gives the same results as the notebook.

Another point we discussed in the forum was to have only one verification mechanism. For
this reason RC and VO use the same tool, so there is only one detector in the whole
system. The tool is not an agent by itself: it becomes part of the MAS through the agents
that decide when to call it and how to read its results.
)
% Needs: tikz, graphicx, booktabs, tabularx, array (already in the preamble).
% =================================================
\DeclareRobustCommand{\pfull}{\tikz[baseline=-0.6ex]\fill (0,0) circle (0.8ex);}
\DeclareRobustCommand{\phalf}{\tikz[baseline=-0.6ex]{\draw (0,0) circle (0.8ex);
  \fill (0,-0.8ex) arc (-90:90:0.8ex) -- cycle;}}
\DeclareRobustCommand{\pnone}{\tikz[baseline=-0.6ex]\draw (0,0) circle (0.8ex);}
\providecommand{\rot}[1]{\rotatebox{75}{#1}}

\section{Agent Taxonomy}
\label{sec:taxonomy}

In this section we analyse the agents of Section~3 with the theory seen in the course:
agent properties, agent architectures and types of agents by purpose
\cite{lecture2, lecture3, wooldridge2009}. To keep the tables short we use these
abbreviations: Flight Data Assistant (FDA), Trajectory Planner (TP), Weather Data
Collector (WDC), Weather Hazard Analyst (WHA), Radar Controller (RC), Conflict Analyst
(CA), Executive Controller (EC), Verification Officer (VO), ATC Supervisor (SUP) and ATC
Report Officer (RO).

% -------------------------------------------------
\subsection{Agent Properties}
\label{sec:properties}

Table~\ref{tab:properties} shows which of the ten properties of Lecture~3 each agent has.
When we filled it in, we saw that three features of our design explain most of the
results. First, the system runs only once and then stops. Second, all the calculations
are done by deterministic tools. Third, the agents exchange information through
intermediate files with a fixed structure (Pydantic models).

\begin{table}[H]
\centering
\small
\caption{Agent properties. \pfull{} = yes, \phalf{} = partially, \pnone{} = no.}
\label{tab:properties}
\begin{tabular}{@{}l*{10}{c}@{}}
\toprule
\textbf{Agent} & \rot{Flexibility} & \rot{Reactivity} & \rot{Proactiveness} & \rot{Social ability}
  & \rot{Rationality} & \rot{Reasoning} & \rot{Learning} & \rot{Autonomy}
  & \rot{Temporal continuity} & \rot{Mobility} \\
\midrule
FDA & \phalf & \phalf & \pnone & \phalf & \pfull & \pnone & \pnone & \phalf & \pnone & \pnone \\
TP  & \pnone & \phalf & \pnone & \phalf & \pfull & \pnone & \pnone & \phalf & \pnone & \pnone \\
WDC & \pnone & \phalf & \phalf & \phalf & \pfull & \pnone & \pnone & \phalf & \pnone & \pnone \\
WHA & \phalf & \phalf & \phalf & \phalf & \pfull & \phalf & \pnone & \phalf & \pnone & \pnone \\
RC  & \pnone & \phalf & \pnone & \phalf & \pfull & \pnone & \pnone & \phalf & \pnone & \pnone \\
CA  & \phalf & \phalf & \pfull & \phalf & \pfull & \pfull & \pnone & \phalf & \pnone & \pnone \\
EC  & \pfull & \pfull & \pfull & \pfull & \pfull & \pfull & \phalf & \pfull & \pnone & \pnone \\
VO  & \pnone & \pfull & \pnone & \phalf & \pfull & \pnone & \pnone & \phalf & \pnone & \pnone \\
SUP & \pfull & \pfull & \pfull & \pfull & \pfull & \pfull & \pnone & \pfull & \pnone & \pnone \\
RO  & \phalf & \phalf & \phalf & \phalf & \pfull & \phalf & \pnone & \phalf & \pnone & \pnone \\
\bottomrule
\end{tabular}
\end{table}

The data and detection agents (FDA, TP, WDC, RC, VO) have few properties, because they only
apply a fixed procedure. On the other hand, EC and SUP have almost all of them, since they
deal with the open part of the problem: deciding how to solve each conflict.

% -------------------------------------------------
\subsection{Justification of Non-obvious Property Classifications}
\label{sec:properties-justification}

\textbf{Reactivity.} Our environment does not change while the system is running
(Section~2), so in the usual sense \cite{lecture3} our agents would not be very reactive.
We decided to look instead at the changes in the information that each agent receives
from the others. Most agents get their input only once, so we marked them as partially
reactive. EC, VO and SUP are different: in each iteration of the verification loop they
receive new flight plans, a rejection or new feedback, and they have to adapt to it.

\textbf{Proactiveness.} CA does more than pass the list of conflicts on. It decides which
ones are more urgent and gives the Tactical crew an order to work on them. RO is
partially proactive, because it points out the critical cases to the human controller
without being asked.

\textbf{Social ability.} Our agents do not use an agent communication language like
FIPA-ACL. They write and read intermediate files, which works in a similar way to a
blackboard system \cite{lecture4}, and the Pydantic models act as a shared ontology. The
only real two-way dialogue happens between EC and SUP: EC proposes a manoeuvre and SUP
accepts or rejects it, which is close to the propose, accept-proposal and reject-proposal
acts of FIPA \cite{fipa_acl}.

\textbf{Rationality.} A rational agent acts towards its goals as far as its beliefs allow
\cite{lecture3}. The problem is that a language model can have wrong beliefs, for example a
wrong distance between two aircraft. That is why all the numbers in our system come from
deterministic tools, and every manoeuvre is checked by the Safety Validation crew before
it reaches the report.

\textbf{Learning.} No agent learns over time: each execution is independent and nothing is
stored between runs. We only marked EC as partial, because inside one run it uses the
feedback of SUP to improve its next proposal. Saving past conflicts and the solutions that
were accepted (case-based reasoning) could be a future improvement.

\textbf{Autonomy.} In the course we saw that the user can limit the autonomy of an agent,
and that some decisions should not be fully delegated to a machine \cite{lecture3}. In ATC
this is very clear. EC is free to choose how to solve a conflict, but it cannot apply its
decision directly: SUP has to approve it, and conflicts that are still unsolved after two
attempts go to a human controller. In the forum we also agreed that the iteration counter
is kept in the flow state and not by the language model, so the agent cannot decide to
skip the limit.

\textbf{Temporal continuity and mobility.} Agents with temporal continuity are running all
the time, and they are not ``once-only computations'' that read an input, produce an
output and stop \cite{lecture3}. Our system works in this second way, so none of our
agents has this property, although a real ATC system would need it. None of them is mobile
either, since all of them run on the same computer.

% -------------------------------------------------
\subsection{Agent Architectures}
\label{sec:architectures}

We classified the agents as reactive, deliberative or hybrid \cite{lecture2}. Since all
CrewAI agents use a language model internally, the question for us was how each agent
makes its decisions: by following a fixed rule, or by reasoning about its goals.
Table~\ref{tab:classification} shows the result.

\begin{table}[H]
\centering
\small
\caption{Architecture and type of each agent.}
\label{tab:classification}
\begin{tabularx}{\textwidth}{@{}l l l X@{}}
\toprule
\textbf{Agent} & \textbf{Architecture} & \textbf{Type} & \textbf{Comment} \\
\midrule
Flight Data Assistant  & Reactive     & Information   & Combines GraphML and JSON sources \\
Trajectory Planner     & Reactive     & Information   & Builds the trajectories for the other crews \\
Weather Data Collector & Reactive     & Information   & Loads the weather snapshot \\
Weather Hazard Analyst & Hybrid       & Information   & Turns weather data into hazard zones \\
Radar Controller       & Reactive     & Information   & Runs the detection tool \\
Conflict Analyst       & Hybrid       & Collaborative & Prioritises conflicts for Tactical \\
Executive Controller   & Deliberative & Collaborative & BDI-style planning of manoeuvres \\
Verification Officer   & Reactive     & Collaborative & Same detection rule as the RC \\
ATC Supervisor         & Hybrid       & Collaborative & Coordinates with the EC \\
ATC Report Officer     & Deliberative & Interface     & Assists the human controller \\
\bottomrule
\end{tabularx}
\end{table}

\textbf{Reactive agents.} FDA, TP, WDC, RC and VO do not make plans. They just map a
situation to an action ($S \rightarrow A$) \cite{lecture2}. RC and VO are a good example,
because their main rule is ``if two aircraft are closer than 5~NM horizontally and
1000~ft vertically, report a conflict'' \cite{icao_doc4444}. It reminded us of the
obstacle-avoidance rule in the reactive layer of TouringMachines \cite{ferguson1992}. We
want this kind of agent here because the check runs in every iteration, so it has to be
fast and reliable.

\textbf{Deliberative agents.} EC can be described with the BDI model \cite{bratman1987}.
Its \emph{beliefs} are the conflicts, the hazard zones, the aircraft performance and the
feedback of SUP. Its \emph{desire} is to solve all the conflicts with as little delay or
extra distance as possible. The manoeuvres it chooses are its \emph{intentions}, and the
concrete actions (climb, reroute or delay) are its \emph{plans}. When SUP rejects a
manoeuvre, EC drops that intention and looks for another one, in line with the idea that
an agent keeps its intentions only while it believes they can be achieved
\cite{cohen1990}. We also see RO as deliberative, because it has to plan the report from
the results of all the crews.

\textbf{Hybrid agents.} WHA, CA and SUP mix a reactive part (thresholds and tools) with a
deliberative part (interpreting, prioritising or deciding). For SUP we found that the
three layers of InteRRaP \cite{muller1996} describe it quite well. Its behaviour layer
rejects any manoeuvre that creates a new conflict. Its planning layer looks for the
reason and suggests something else. Its cooperative layer sends structured feedback to EC.

If we look at the whole system, it also has a hybrid structure. The tools work as a
reactive layer, the language-model agents as a deliberative layer, and the CrewAI Flow
with its routers decides which one acts, a bit like the control framework of
TouringMachines. The routers are deterministic, so the most critical decisions never
depend on the language model.

% -------------------------------------------------
\subsection{Agent Types by Purpose}
\label{sec:types}

Agents can also be classified by purpose as collaborative, interface, information or
facilitator agents \cite{lecture3, nwana1996}. Our system has several of these types
working together.

\textbf{Collaborative agents.} Air traffic control is one of the examples of collaborative
agents given in the course \cite{lecture3}, and in our design CA, EC, VO and SUP solve the
conflicts together through proposals, checks and feedback. Since all the agents share the
same goal, they are benevolent. Using the cooperation hierarchy of Lecture~5
\cite{doran1997}, we would place our system among the cooperative MAS with explicit
communication between deliberative agents. It also follows the steps of distributed
problem solving \cite{lecture4}. The problem is split into crews at design time, each crew
solves one part, the results are shared through the intermediate files, and the Reporting
crew puts them together in the final report.

\textbf{Information agents.} These agents collect, filter and combine data from different
sources \cite{lecture3}. FDA merges three input files that come in different formats, WDC
loads the weather data, and WHA turns it into hazard zones. TP and RC then give the
decision agents a clear view of the trajectories and the conflicts. In the forum we chose
to load a stored weather file instead of downloading live data, so WDC always works with
the same snapshot and the results can be repeated.

\textbf{Interface agent.} RO is the only agent that works for a person, the controller who
reads the report. We call it an interface agent only in a limited sense, because interface
agents usually learn from their user \cite{lecture3} and ours does not.

\textbf{Facilitators.} There is no facilitator agent in our design. The CrewAI Flow takes
this role: it decides which crew runs next and passes the data between them. We think a
facilitator agent would only make sense if new agents could join the system while it is
running.

% -------------------------------------------------
\subsection{Deterministic Conflict Tool and Agentification}
\label{sec:agentification}

The conflict detection code from the course notebook is not an agent, so we had to decide
how to connect it to the MAS. The course presents three options: a translator, a wrapper,
or rewriting the code as an agent \cite{lecture3, genesereth1994}.

A translator was not a good option, because the notebook calculates the validation errors
but does not return them, and its output has no structure. Rewriting it as an agent would
mean letting the language model calculate the distances, and we did not want that in a
safety-critical system. In the end we chose a \textbf{wrapper}. We moved the code into a
Python module with small changes (Pydantic models and returning the validation errors)
and added a CrewAI tool on top of it \cite{crewai_docs}. The algorithm is the same, and our
tests check that it gives the same results as the notebook.

Another point we discussed in the forum was to have only one verification mechanism. For
this reason RC and VO use the same tool, so there is only one detector in the whole
system. The tool is not an agent by itself: it becomes part of the MAS through the agents
that decide when to call it and how to read its results.
)
% Needs: tikz, graphicx, booktabs, tabularx, array (already in the preamble).
% =================================================
\DeclareRobustCommand{\pfull}{\tikz[baseline=-0.6ex]\fill (0,0) circle (0.8ex);}
\DeclareRobustCommand{\phalf}{\tikz[baseline=-0.6ex]{\draw (0,0) circle (0.8ex);
  \fill (0,-0.8ex) arc (-90:90:0.8ex) -- cycle;}}
\DeclareRobustCommand{\pnone}{\tikz[baseline=-0.6ex]\draw (0,0) circle (0.8ex);}
\providecommand{\rot}[1]{\rotatebox{75}{#1}}

\section{Agent Taxonomy}
\label{sec:taxonomy}

In this section we analyse the agents of Section~3 with the theory seen in the course:
agent properties, agent architectures and types of agents by purpose
\cite{lecture2, lecture3, wooldridge2009}. To keep the tables short we use these
abbreviations: Flight Data Assistant (FDA), Trajectory Planner (TP), Weather Data
Collector (WDC), Weather Hazard Analyst (WHA), Radar Controller (RC), Conflict Analyst
(CA), Executive Controller (EC), Verification Officer (VO), ATC Supervisor (SUP) and ATC
Report Officer (RO).

% -------------------------------------------------
\subsection{Agent Properties}
\label{sec:properties}

Table~\ref{tab:properties} shows which of the ten properties of Lecture~3 each agent has.
When we filled it in, we saw that three features of our design explain most of the
results. First, the system runs only once and then stops. Second, all the calculations
are done by deterministic tools. Third, the agents exchange information through
intermediate files with a fixed structure (Pydantic models).

\begin{table}[H]
\centering
\small
\caption{Agent properties. \pfull{} = yes, \phalf{} = partially, \pnone{} = no.}
\label{tab:properties}
\begin{tabular}{@{}l*{10}{c}@{}}
\toprule
\textbf{Agent} & \rot{Flexibility} & \rot{Reactivity} & \rot{Proactiveness} & \rot{Social ability}
  & \rot{Rationality} & \rot{Reasoning} & \rot{Learning} & \rot{Autonomy}
  & \rot{Temporal continuity} & \rot{Mobility} \\
\midrule
FDA & \phalf & \phalf & \pnone & \phalf & \pfull & \pnone & \pnone & \phalf & \pnone & \pnone \\
TP  & \pnone & \phalf & \pnone & \phalf & \pfull & \pnone & \pnone & \phalf & \pnone & \pnone \\
WDC & \pnone & \phalf & \phalf & \phalf & \pfull & \pnone & \pnone & \phalf & \pnone & \pnone \\
WHA & \phalf & \phalf & \phalf & \phalf & \pfull & \phalf & \pnone & \phalf & \pnone & \pnone \\
RC  & \pnone & \phalf & \pnone & \phalf & \pfull & \pnone & \pnone & \phalf & \pnone & \pnone \\
CA  & \phalf & \phalf & \pfull & \phalf & \pfull & \pfull & \pnone & \phalf & \pnone & \pnone \\
EC  & \pfull & \pfull & \pfull & \pfull & \pfull & \pfull & \phalf & \pfull & \pnone & \pnone \\
VO  & \pnone & \pfull & \pnone & \phalf & \pfull & \pnone & \pnone & \phalf & \pnone & \pnone \\
SUP & \pfull & \pfull & \pfull & \pfull & \pfull & \pfull & \pnone & \pfull & \pnone & \pnone \\
RO  & \phalf & \phalf & \phalf & \phalf & \pfull & \phalf & \pnone & \phalf & \pnone & \pnone \\
\bottomrule
\end{tabular}
\end{table}

The data and detection agents (FDA, TP, WDC, RC, VO) have few properties, because they only
apply a fixed procedure. On the other hand, EC and SUP have almost all of them, since they
deal with the open part of the problem: deciding how to solve each conflict.

% -------------------------------------------------
\subsection{Justification of Non-obvious Property Classifications}
\label{sec:properties-justification}

\textbf{Reactivity.} Our environment does not change while the system is running
(Section~2), so in the usual sense \cite{lecture3} our agents would not be very reactive.
We decided to look instead at the changes in the information that each agent receives
from the others. Most agents get their input only once, so we marked them as partially
reactive. EC, VO and SUP are different: in each iteration of the verification loop they
receive new flight plans, a rejection or new feedback, and they have to adapt to it.

\textbf{Proactiveness.} CA does more than pass the list of conflicts on. It decides which
ones are more urgent and gives the Tactical crew an order to work on them. RO is
partially proactive, because it points out the critical cases to the human controller
without being asked.

\textbf{Social ability.} Our agents do not use an agent communication language like
FIPA-ACL. They write and read intermediate files, which works in a similar way to a
blackboard system \cite{lecture4}, and the Pydantic models act as a shared ontology. The
only real two-way dialogue happens between EC and SUP: EC proposes a manoeuvre and SUP
accepts or rejects it, which is close to the propose, accept-proposal and reject-proposal
acts of FIPA \cite{fipa_acl}.

\textbf{Rationality.} A rational agent acts towards its goals as far as its beliefs allow
\cite{lecture3}. The problem is that a language model can have wrong beliefs, for example a
wrong distance between two aircraft. That is why all the numbers in our system come from
deterministic tools, and every manoeuvre is checked by the Safety Validation crew before
it reaches the report.

\textbf{Learning.} No agent learns over time: each execution is independent and nothing is
stored between runs. We only marked EC as partial, because inside one run it uses the
feedback of SUP to improve its next proposal. Saving past conflicts and the solutions that
were accepted (case-based reasoning) could be a future improvement.

\textbf{Autonomy.} In the course we saw that the user can limit the autonomy of an agent,
and that some decisions should not be fully delegated to a machine \cite{lecture3}. In ATC
this is very clear. EC is free to choose how to solve a conflict, but it cannot apply its
decision directly: SUP has to approve it, and conflicts that are still unsolved after two
attempts go to a human controller. In the forum we also agreed that the iteration counter
is kept in the flow state and not by the language model, so the agent cannot decide to
skip the limit.

\textbf{Temporal continuity and mobility.} Agents with temporal continuity are running all
the time, and they are not ``once-only computations'' that read an input, produce an
output and stop \cite{lecture3}. Our system works in this second way, so none of our
agents has this property, although a real ATC system would need it. None of them is mobile
either, since all of them run on the same computer.

% -------------------------------------------------
\subsection{Agent Architectures}
\label{sec:architectures}

We classified the agents as reactive, deliberative or hybrid \cite{lecture2}. Since all
CrewAI agents use a language model internally, the question for us was how each agent
makes its decisions: by following a fixed rule, or by reasoning about its goals.
Table~\ref{tab:classification} shows the result.

\begin{table}[H]
\centering
\small
\caption{Architecture and type of each agent.}
\label{tab:classification}
\begin{tabularx}{\textwidth}{@{}l l l X@{}}
\toprule
\textbf{Agent} & \textbf{Architecture} & \textbf{Type} & \textbf{Comment} \\
\midrule
Flight Data Assistant  & Reactive     & Information   & Combines GraphML and JSON sources \\
Trajectory Planner     & Reactive     & Information   & Builds the trajectories for the other crews \\
Weather Data Collector & Reactive     & Information   & Loads the weather snapshot \\
Weather Hazard Analyst & Hybrid       & Information   & Turns weather data into hazard zones \\
Radar Controller       & Reactive     & Information   & Runs the detection tool \\
Conflict Analyst       & Hybrid       & Collaborative & Prioritises conflicts for Tactical \\
Executive Controller   & Deliberative & Collaborative & BDI-style planning of manoeuvres \\
Verification Officer   & Reactive     & Collaborative & Same detection rule as the RC \\
ATC Supervisor         & Hybrid       & Collaborative & Coordinates with the EC \\
ATC Report Officer     & Deliberative & Interface     & Assists the human controller \\
\bottomrule
\end{tabularx}
\end{table}

\textbf{Reactive agents.} FDA, TP, WDC, RC and VO do not make plans. They just map a
situation to an action ($S \rightarrow A$) \cite{lecture2}. RC and VO are a good example,
because their main rule is ``if two aircraft are closer than 5~NM horizontally and
1000~ft vertically, report a conflict'' \cite{icao_doc4444}. It reminded us of the
obstacle-avoidance rule in the reactive layer of TouringMachines \cite{ferguson1992}. We
want this kind of agent here because the check runs in every iteration, so it has to be
fast and reliable.

\textbf{Deliberative agents.} EC can be described with the BDI model \cite{bratman1987}.
Its \emph{beliefs} are the conflicts, the hazard zones, the aircraft performance and the
feedback of SUP. Its \emph{desire} is to solve all the conflicts with as little delay or
extra distance as possible. The manoeuvres it chooses are its \emph{intentions}, and the
concrete actions (climb, reroute or delay) are its \emph{plans}. When SUP rejects a
manoeuvre, EC drops that intention and looks for another one, in line with the idea that
an agent keeps its intentions only while it believes they can be achieved
\cite{cohen1990}. We also see RO as deliberative, because it has to plan the report from
the results of all the crews.

\textbf{Hybrid agents.} WHA, CA and SUP mix a reactive part (thresholds and tools) with a
deliberative part (interpreting, prioritising or deciding). For SUP we found that the
three layers of InteRRaP \cite{muller1996} describe it quite well. Its behaviour layer
rejects any manoeuvre that creates a new conflict. Its planning layer looks for the
reason and suggests something else. Its cooperative layer sends structured feedback to EC.

If we look at the whole system, it also has a hybrid structure. The tools work as a
reactive layer, the language-model agents as a deliberative layer, and the CrewAI Flow
with its routers decides which one acts, a bit like the control framework of
TouringMachines. The routers are deterministic, so the most critical decisions never
depend on the language model.

% -------------------------------------------------
\subsection{Agent Types by Purpose}
\label{sec:types}

Agents can also be classified by purpose as collaborative, interface, information or
facilitator agents \cite{lecture3, nwana1996}. Our system has several of these types
working together.

\textbf{Collaborative agents.} Air traffic control is one of the examples of collaborative
agents given in the course \cite{lecture3}, and in our design CA, EC, VO and SUP solve the
conflicts together through proposals, checks and feedback. Since all the agents share the
same goal, they are benevolent. Using the cooperation hierarchy of Lecture~5
\cite{doran1997}, we would place our system among the cooperative MAS with explicit
communication between deliberative agents. It also follows the steps of distributed
problem solving \cite{lecture4}. The problem is split into crews at design time, each crew
solves one part, the results are shared through the intermediate files, and the Reporting
crew puts them together in the final report.

\textbf{Information agents.} These agents collect, filter and combine data from different
sources \cite{lecture3}. FDA merges three input files that come in different formats, WDC
loads the weather data, and WHA turns it into hazard zones. TP and RC then give the
decision agents a clear view of the trajectories and the conflicts. In the forum we chose
to load a stored weather file instead of downloading live data, so WDC always works with
the same snapshot and the results can be repeated.

\textbf{Interface agent.} RO is the only agent that works for a person, the controller who
reads the report. We call it an interface agent only in a limited sense, because interface
agents usually learn from their user \cite{lecture3} and ours does not.

\textbf{Facilitators.} There is no facilitator agent in our design. The CrewAI Flow takes
this role: it decides which crew runs next and passes the data between them. We think a
facilitator agent would only make sense if new agents could join the system while it is
running.

% -------------------------------------------------
\subsection{Deterministic Conflict Tool and Agentification}
\label{sec:agentification}

The conflict detection code from the course notebook is not an agent, so we had to decide
how to connect it to the MAS. The course presents three options: a translator, a wrapper,
or rewriting the code as an agent \cite{lecture3, genesereth1994}.

A translator was not a good option, because the notebook calculates the validation errors
but does not return them, and its output has no structure. Rewriting it as an agent would
mean letting the language model calculate the distances, and we did not want that in a
safety-critical system. In the end we chose a \textbf{wrapper}. We moved the code into a
Python module with small changes (Pydantic models and returning the validation errors)
and added a CrewAI tool on top of it \cite{crewai_docs}. The algorithm is the same, and our
tests check that it gives the same results as the notebook.

Another point we discussed in the forum was to have only one verification mechanism. For
this reason RC and VO use the same tool, so there is only one detector in the whole
system. The tool is not an agent by itself: it becomes part of the MAS through the agents
that decide when to call it and how to read its results.
)
% Needs: tikz, graphicx, booktabs, tabularx, array (already in the preamble).
% =================================================
\DeclareRobustCommand{\pfull}{\tikz[baseline=-0.6ex]\fill (0,0) circle (0.8ex);}
\DeclareRobustCommand{\phalf}{\tikz[baseline=-0.6ex]{\draw (0,0) circle (0.8ex);
  \fill (0,-0.8ex) arc (-90:90:0.8ex) -- cycle;}}
\DeclareRobustCommand{\pnone}{\tikz[baseline=-0.6ex]\draw (0,0) circle (0.8ex);}
\providecommand{\rot}[1]{\rotatebox{75}{#1}}

\section{Agent Taxonomy}
\label{sec:taxonomy}

In this section we analyse the agents of Section~3 with the theory seen in the course:
agent properties, agent architectures and types of agents by purpose
\cite{lecture2, lecture3, wooldridge2009}. To keep the tables short we use these
abbreviations: Flight Data Assistant (FDA), Trajectory Planner (TP), Weather Data
Collector (WDC), Weather Hazard Analyst (WHA), Radar Controller (RC), Conflict Analyst
(CA), Executive Controller (EC), Verification Officer (VO), ATC Supervisor (SUP) and ATC
Report Officer (RO).

% -------------------------------------------------
\subsection{Agent Properties}
\label{sec:properties}

Table~\ref{tab:properties} shows which of the ten properties of Lecture~3 each agent has.
When we filled it in, we saw that three features of our design explain most of the
results. First, the system runs only once and then stops. Second, all the calculations
are done by deterministic tools. Third, the agents exchange information through
intermediate files with a fixed structure (Pydantic models).

\begin{table}[H]
\centering
\small
\caption{Agent properties. \pfull{} = yes, \phalf{} = partially, \pnone{} = no.}
\label{tab:properties}
\begin{tabular}{@{}l*{10}{c}@{}}
\toprule
\textbf{Agent} & \rot{Flexibility} & \rot{Reactivity} & \rot{Proactiveness} & \rot{Social ability}
  & \rot{Rationality} & \rot{Reasoning} & \rot{Learning} & \rot{Autonomy}
  & \rot{Temporal continuity} & \rot{Mobility} \\
\midrule
FDA & \phalf & \phalf & \pnone & \phalf & \pfull & \pnone & \pnone & \phalf & \pnone & \pnone \\
TP  & \pnone & \phalf & \pnone & \phalf & \pfull & \pnone & \pnone & \phalf & \pnone & \pnone \\
WDC & \pnone & \phalf & \phalf & \phalf & \pfull & \pnone & \pnone & \phalf & \pnone & \pnone \\
WHA & \phalf & \phalf & \phalf & \phalf & \pfull & \phalf & \pnone & \phalf & \pnone & \pnone \\
RC  & \pnone & \phalf & \pnone & \phalf & \pfull & \pnone & \pnone & \phalf & \pnone & \pnone \\
CA  & \phalf & \phalf & \pfull & \phalf & \pfull & \pfull & \pnone & \phalf & \pnone & \pnone \\
EC  & \pfull & \pfull & \pfull & \pfull & \pfull & \pfull & \phalf & \pfull & \pnone & \pnone \\
VO  & \pnone & \pfull & \pnone & \phalf & \pfull & \pnone & \pnone & \phalf & \pnone & \pnone \\
SUP & \pfull & \pfull & \pfull & \pfull & \pfull & \pfull & \pnone & \pfull & \pnone & \pnone \\
RO  & \phalf & \phalf & \phalf & \phalf & \pfull & \phalf & \pnone & \phalf & \pnone & \pnone \\
\bottomrule
\end{tabular}
\end{table}

The data and detection agents (FDA, TP, WDC, RC, VO) have few properties, because they only
apply a fixed procedure. On the other hand, EC and SUP have almost all of them, since they
deal with the open part of the problem: deciding how to solve each conflict.

% -------------------------------------------------
\subsection{Justification of Non-obvious Property Classifications}
\label{sec:properties-justification}

\textbf{Reactivity.} Our environment does not change while the system is running
(Section~2), so in the usual sense \cite{lecture3} our agents would not be very reactive.
We decided to look instead at the changes in the information that each agent receives
from the others. Most agents get their input only once, so we marked them as partially
reactive. EC, VO and SUP are different: in each iteration of the verification loop they
receive new flight plans, a rejection or new feedback, and they have to adapt to it.

\textbf{Proactiveness.} CA does more than pass the list of conflicts on. It decides which
ones are more urgent and gives the Tactical crew an order to work on them. RO is
partially proactive, because it points out the critical cases to the human controller
without being asked.

\textbf{Social ability.} Our agents do not use an agent communication language like
FIPA-ACL. They write and read intermediate files, which works in a similar way to a
blackboard system \cite{lecture4}, and the Pydantic models act as a shared ontology. The
only real two-way dialogue happens between EC and SUP: EC proposes a manoeuvre and SUP
accepts or rejects it, which is close to the propose, accept-proposal and reject-proposal
acts of FIPA \cite{fipa_acl}.

\textbf{Rationality.} A rational agent acts towards its goals as far as its beliefs allow
\cite{lecture3}. The problem is that a language model can have wrong beliefs, for example a
wrong distance between two aircraft. That is why all the numbers in our system come from
deterministic tools, and every manoeuvre is checked by the Safety Validation crew before
it reaches the report.

\textbf{Learning.} No agent learns over time: each execution is independent and nothing is
stored between runs. We only marked EC as partial, because inside one run it uses the
feedback of SUP to improve its next proposal. Saving past conflicts and the solutions that
were accepted (case-based reasoning) could be a future improvement.

\textbf{Autonomy.} In the course we saw that the user can limit the autonomy of an agent,
and that some decisions should not be fully delegated to a machine \cite{lecture3}. In ATC
this is very clear. EC is free to choose how to solve a conflict, but it cannot apply its
decision directly: SUP has to approve it, and conflicts that are still unsolved after two
attempts go to a human controller. In the forum we also agreed that the iteration counter
is kept in the flow state and not by the language model, so the agent cannot decide to
skip the limit.

\textbf{Temporal continuity and mobility.} Agents with temporal continuity are running all
the time, and they are not ``once-only computations'' that read an input, produce an
output and stop \cite{lecture3}. Our system works in this second way, so none of our
agents has this property, although a real ATC system would need it. None of them is mobile
either, since all of them run on the same computer.

% -------------------------------------------------
\subsection{Agent Architectures}
\label{sec:architectures}

We classified the agents as reactive, deliberative or hybrid \cite{lecture2}. Since all
CrewAI agents use a language model internally, the question for us was how each agent
makes its decisions: by following a fixed rule, or by reasoning about its goals.
Table~\ref{tab:classification} shows the result.

\begin{table}[H]
\centering
\small
\caption{Architecture and type of each agent.}
\label{tab:classification}
\begin{tabularx}{\textwidth}{@{}l l l X@{}}
\toprule
\textbf{Agent} & \textbf{Architecture} & \textbf{Type} & \textbf{Comment} \\
\midrule
Flight Data Assistant  & Reactive     & Information   & Combines GraphML and JSON sources \\
Trajectory Planner     & Reactive     & Information   & Builds the trajectories for the other crews \\
Weather Data Collector & Reactive     & Information   & Loads the weather snapshot \\
Weather Hazard Analyst & Hybrid       & Information   & Turns weather data into hazard zones \\
Radar Controller       & Reactive     & Information   & Runs the detection tool \\
Conflict Analyst       & Hybrid       & Collaborative & Prioritises conflicts for Tactical \\
Executive Controller   & Deliberative & Collaborative & BDI-style planning of manoeuvres \\
Verification Officer   & Reactive     & Collaborative & Same detection rule as the RC \\
ATC Supervisor         & Hybrid       & Collaborative & Coordinates with the EC \\
ATC Report Officer     & Deliberative & Interface     & Assists the human controller \\
\bottomrule
\end{tabularx}
\end{table}

\textbf{Reactive agents.} FDA, TP, WDC, RC and VO do not make plans. They just map a
situation to an action ($S \rightarrow A$) \cite{lecture2}. RC and VO are a good example,
because their main rule is ``if two aircraft are closer than 5~NM horizontally and
1000~ft vertically, report a conflict'' \cite{icao_doc4444}. It reminded us of the
obstacle-avoidance rule in the reactive layer of TouringMachines \cite{ferguson1992}. We
want this kind of agent here because the check runs in every iteration, so it has to be
fast and reliable.

\textbf{Deliberative agents.} EC can be described with the BDI model \cite{bratman1987}.
Its \emph{beliefs} are the conflicts, the hazard zones, the aircraft performance and the
feedback of SUP. Its \emph{desire} is to solve all the conflicts with as little delay or
extra distance as possible. The manoeuvres it chooses are its \emph{intentions}, and the
concrete actions (climb, reroute or delay) are its \emph{plans}. When SUP rejects a
manoeuvre, EC drops that intention and looks for another one, in line with the idea that
an agent keeps its intentions only while it believes they can be achieved
\cite{cohen1990}. We also see RO as deliberative, because it has to plan the report from
the results of all the crews.

\textbf{Hybrid agents.} WHA, CA and SUP mix a reactive part (thresholds and tools) with a
deliberative part (interpreting, prioritising or deciding). For SUP we found that the
three layers of InteRRaP \cite{muller1996} describe it quite well. Its behaviour layer
rejects any manoeuvre that creates a new conflict. Its planning layer looks for the
reason and suggests something else. Its cooperative layer sends structured feedback to EC.

If we look at the whole system, it also has a hybrid structure. The tools work as a
reactive layer, the language-model agents as a deliberative layer, and the CrewAI Flow
with its routers decides which one acts, a bit like the control framework of
TouringMachines. The routers are deterministic, so the most critical decisions never
depend on the language model.

% -------------------------------------------------
\subsection{Agent Types by Purpose}
\label{sec:types}

Agents can also be classified by purpose as collaborative, interface, information or
facilitator agents \cite{lecture3, nwana1996}. Our system has several of these types
working together.

\textbf{Collaborative agents.} Air traffic control is one of the examples of collaborative
agents given in the course \cite{lecture3}, and in our design CA, EC, VO and SUP solve the
conflicts together through proposals, checks and feedback. Since all the agents share the
same goal, they are benevolent. Using the cooperation hierarchy of Lecture~5
\cite{doran1997}, we would place our system among the cooperative MAS with explicit
communication between deliberative agents. It also follows the steps of distributed
problem solving \cite{lecture4}. The problem is split into crews at design time, each crew
solves one part, the results are shared through the intermediate files, and the Reporting
crew puts them together in the final report.

\textbf{Information agents.} These agents collect, filter and combine data from different
sources \cite{lecture3}. FDA merges three input files that come in different formats, WDC
loads the weather data, and WHA turns it into hazard zones. TP and RC then give the
decision agents a clear view of the trajectories and the conflicts. In the forum we chose
to load a stored weather file instead of downloading live data, so WDC always works with
the same snapshot and the results can be repeated.

\textbf{Interface agent.} RO is the only agent that works for a person, the controller who
reads the report. We call it an interface agent only in a limited sense, because interface
agents usually learn from their user \cite{lecture3} and ours does not.

\textbf{Facilitators.} There is no facilitator agent in our design. The CrewAI Flow takes
this role: it decides which crew runs next and passes the data between them. We think a
facilitator agent would only make sense if new agents could join the system while it is
running.

% -------------------------------------------------
\subsection{Deterministic Conflict Tool and Agentification}
\label{sec:agentification}

The conflict detection code from the course notebook is not an agent, so we had to decide
how to connect it to the MAS. The course presents three options: a translator, a wrapper,
or rewriting the code as an agent \cite{lecture3, genesereth1994}.

A translator was not a good option, because the notebook calculates the validation errors
but does not return them, and its output has no structure. Rewriting it as an agent would
mean letting the language model calculate the distances, and we did not want that in a
safety-critical system. In the end we chose a \textbf{wrapper}. We moved the code into a
Python module with small changes (Pydantic models and returning the validation errors)
and added a CrewAI tool on top of it \cite{crewai_docs}. The algorithm is the same, and our
tests check that it gives the same results as the notebook.

Another point we discussed in the forum was to have only one verification mechanism. For
this reason RC and VO use the same tool, so there is only one detector in the whole
system. The tool is not an agent by itself: it becomes part of the MAS through the agents
that decide when to call it and how to read its results.
